% Technical Background, section 3. Pointed back to by: methodology/03_centerline_tree.tex (section
% body); 06_bifurcation_angle.tex and 07_daughter_ratio.tex (P4b, radius); 04_artery_labeling.tex
% (subsection 3.1).
% Model: Aida 3.3 (FEM, section body) and 3.3.1 (the subsection one step deeper).
% Cite: antiga2008vmtk (minimal-path / Voronoi centerlines), hampe2024gnn (tool, coronary labelling
%   with a GNN).
% CITE WANTED: a skeletonisation (topological thinning) source, e.g. Lee et al. 1994.
% CITE WANTED: the in-house tracer's publication, if one exists.
% CITE WANTED: a GNN foundation, Gilmer et al. 2017 (message passing) or Kipf and Welling 2017.
% Medis is NOT described here: it is a commercial product that delivered reference labels, so it
%   belongs in the Data chapter and the Methodology (methodology.tex, Artery Labeling).
% FIGURE WANTED: none; the Methodology may show a mask, its skeleton and the graph instead.

\section{Centerline Extraction and the Coronary Tree}
\label{sec:tech-centerline}
% P1 definition opener: a centerline is the curve running through the middle of a vessel lumen.
% P2 skeletonisation thins a binary segmentation to a one-voxel-wide skeleton; it inherits every
%    error of the segmentation.
% P3 tracking or minimal-path methods follow the vessel in the image or a distance map from seed
%    points (antiga2008vmtk); they need seeds and can leave the vessel at low contrast.
% P4 PLACEHOLDER, in-house tracer: which of the two families it belongs to, its input (image or
%    segmentation) and what it adds (e.g. radius estimate, branch handling). Settings, version and
%    "used without modification" go in the Methodology.
% P4b radius: each centerline point carries the radius of the largest sphere inscribed in the lumen
%    there (antiga2008vmtk), from a distance transform of the mask. Near a junction that sphere
%    spans the merged lumen of parent and daughters, so radius and direction there describe the
%    junction, not either vessel. Fact only; the skip and window that resolve it are Methodology.
% P5 a centerline network as a graph: junctions and termini as nodes, segments as edges; rooting at
%    the ostium and orienting away from it gives parent, children and generation.
% P6 HAZARD CLOSER: a break in the lumen becomes a disconnected fragment and a false junction a
%    loop, so every feature downstream depends on the extraction being topologically correct.

\subsection{Graph Neural Networks}
\label{sec:tech-gnn}
% P1 subsection-as-answer opener: a convolution (Section ref sec:tech-unet) needs a regular grid,
%    while a tree has none; a GNN learns on the graph itself.
% P2 message passing: each node updates its features from its neighbours'; k layers reach k hops.
%    One equation for the update, with a "where" sentence.
% P3 why it suits artery labelling: an artery's name depends on its position in the tree, not only
%    on its own shape; the output per segment is a class.
% P4 tool: hampe2024gnn applied to coronary centerline trees.
